\section{Security in virtual machines}
Virtualisation technologies fundamentally redefine the way we see cybersecurity, bringing both advantages for system protection and also a completly new attack surface that did not exist in traditional physical Infrastructure \cite{vm-security.pdf,vm-attack-vectors}. Some examples of advantages brought by virtualisation are: \textbf{Strong isolation} which makes each virtual machine to run in an eviroment completly separated from one another. Even if a guest operating system is initially compromised by an attacker they can't directly access the data or memory of the other virtual machines on the same host \cite{vm-security.pdf}. \textbf{Snapshotting and rapid restoration} gives the administrators the ability to quickly revert the virtual machine's state to a previous one or a clean state in a matter of seconds therefore minimizing the downtime in the event of a ransomware infection or a major system failure. Lasly another advantages brought by the virtual machine is the \textbf{safe analysis and sandboxing}. Virtual machines provide a controlled lab enviroment which is ideal for running and analysing malitions software by detonating it and analysing the behaviour in a secured space. The hypervisor can monitor the behaviour of the application while being completly separated from the virtual machine, this approach is called out-of-band monitoring \cite{vm-security.pdf}. 

\subsection{Attack vectors}
Beside all the good things the virtualisation process adds to the development of the IT Infrastructure it also adds certain vulnerabilities and attack vectors, due to the way virtual machines share physical resources with the host machine and also due to the complexity of the hypervisor \cite{vm-security.pdf}. A specialist should be aware of the risks and attack vectors the usage of virtual machines involves in order to be able to assure their security.

\begin{enumerate}
    \item \textbf{VM Escape}: This attack represents the most severe security incident in resource virtualisation \cite{vm-attack-vectors,vm-security.pdf}. An attacker who succesfully discovered a privilege excalation exploit and gain root access to the virtual machine could further exploit a vulnerability within the hypervisor's hardware emulation code, such as bugs in the emulation of USB adapters or video and network controllers, in order to execute arbitrary code outside of the virtual machine space. Therefore, the attacker can directly obtain full ring 0 privileges on the host operating sytem which allows them to control, inspect or destroy all the other virtual machines \cite{vm-attack-vectors,vm-security.pdf}. 

    \item \textbf{Host-Guest communication channel vulnerabilities}: In order to enhance the user experience, hypervisors provide fast an reliable communication channels between host and guest machines through the use of features like shared clipboard and shared folders. While these channels are designed as a quality of life for the specialist they can also be used by attackers or malware to leak sensitive information outside the secured perimeter or to directly infect the host system files \cite{vm-security.pdf}.
    
    \item \textbf{Denial of Service via resource starvation}: In a virtualised enviroment where hardware resources, such as CPU, RAM, I/Os, are shared a compromised virtual machine or a virtual machine controlled by an attacker could execute an attack that starvates all the other virtual machines of resources to run properly. An example of such attack is the 'noisy neighbor' which is characterized by the creation of infinite loops or the generation of massive I/O transactions that can consume the entire capacity of the physical server resources and starvation of the other legimite VMs or more likely crashing them \cite{vm-security.pdf}.  

    \item \textbf{Packet Sniffing and inter-VM ARP Poisoning}: A compromised virtual machine that gives the attacker root access to the machine could be used to perform ARP poisoning attacks or MAC Flooding attacks in order to intercept and analyse all the traffic performed by the virtual machines that make use of the same virtual switch \cite{vm-security.pdf,vm-attack-vectors}.


    \end{enumerate}

\subsection{Modern security mitigation techniques}
In order to mitigate those vulnerabilities and attack vectors, modern enterprise hypervisors make use of physical chip controlled security technologies. One example of such technology is \textbf{Micro-Hypervisor security}. This method is achieved by reducing the attack surface by stripping off the hypervisor's kernel the complex hardware emulators and by isolating every emulation driver into its own process without elevated privileges. Another example is the usage of \textbf{AMD SEV and Intel SGX}. This hardware assisted technologies make use of cryptographic keys managed at physycal chip level to encrypt the data of the virtual machines that is also separated in RAM. Even if an attacker manages to compromise the host system or somehow manages to gain ring 0 privileges the RAM memory of the guest virtual machine will still be encrypted and completly innacesible to the attacker \cite{vm-attack-vectors,vm-security.pdf}. 